%! AP Statistics — Unit 5, Lesson 5.4 — Guided Notes
\documentclass[10pt]{article}
\usepackage{apstats-article}
\usepackage{apstats-boxes}

\begin{document}

\pageheader{Unit 5, Lesson 5.4}{Guided Notes}
\namedateperiod

\vspace{0.1in}

\begin{objectivebox}
\textbf{Learning Targets}
\begin{enumerate}
  \item I can define bias and explain what makes an estimator unbiased.
  \item I can identify $\bar{x}$, $\hat{p}$, and $s^2$ as unbiased estimators.
  \item I can explain how sample size affects the variability of an estimator.
\end{enumerate}
\end{objectivebox}

\begin{vocabbox}
\begin{description}
  \item[\termblanklong{Point estimator}] a \blank{1.5in} (formula) used to estimate a population parameter.
  \item[\termblanklong{Point estimate}] the \blank{1.5in} (number) computed from sample data.
  \item[\termblanklong{Unbiased estimator}] a statistic whose sampling distribution is centered at the \blank{1.2in}: $E(\hat{\theta}) = \theta$.
  \item[\termblanklong{Bias}] when $E(\hat{\theta}) \ne \theta$; the estimator \emph{systematically} over- or under-estimates the parameter.
\end{description}
\end{vocabbox}

\begin{notesbox}
\textbf{Unbiased Estimators} (UNC-3.I)

\begin{center}
\begin{tabular}{@{} l c c @{}}
\hline
\textbf{Statistic} & \textbf{Estimates} & \textbf{Unbiased?} \\
\hline
Sample mean $\bar{x}$ & $\mu$ & \blank{0.5in} \\
Sample proportion $\hat{p}$ & $p$ & \blank{0.5in} \\
Sample variance $s^2$ (divides by $n-1$) & $\sigma^2$ & \blank{0.5in} \\
Sample range (max $-$ min) & population range & \blank{0.5in} \\
\hline
\end{tabular}
\end{center}

\medskip
Why divide by $n-1$ for variance? \writelines{2}
\end{notesbox}

\begin{notesbox}
\textbf{Variability of Estimators} (UNC-3.J)

\medskip
\begin{tabular}{@{} p{0.45\linewidth} p{0.45\linewidth} @{}}
Larger $n$ $\Rightarrow$ \blank{1.2in} spread & Smaller $n$ $\Rightarrow$ \blank{1.2in} spread
\end{tabular}

\medskip
For the sampling distribution of $\bar{x}$: $\sigma_{\bar{x}} = \sigma/\sqrt{n}$.
\begin{itemize}
  \item Doubling $n$ multiplies $\sigma_{\bar{x}}$ by \blank{0.5in}.
  \item Quadrupling $n$ multiplies $\sigma_{\bar{x}}$ by \blank{0.5in}.
\end{itemize}

\textit{Worked example.} A population has $\sigma = 24$.
\begin{itemize}
  \item $n = 16$: $\sigma_{\bar{x}} = $ \blank{0.5in} \hspace{0.5cm} $n = 64$: $\sigma_{\bar{x}} = $ \blank{0.5in}
\end{itemize}
\end{notesbox}

\begin{notesbox}
\textbf{Bias vs.\ Variability — The Dart Board}

\medskip
\begin{tabular}{@{} p{0.2\linewidth} p{0.7\linewidth} @{}}
Low bias, low variability & \blank{3.5in} \\[6pt]
Low bias, high variability & \blank{3.5in} \\[6pt]
High bias, low variability & \blank{3.5in} \\
\end{tabular}

\medskip
Which is most desirable for a statistical estimator? \writelines{1}
\end{notesbox}

\begin{practicebox}
\textbf{Quick Check}
\begin{enumerate}
  \item A researcher records the \emph{median} salary in many random samples of size 40.
        The distribution of sample medians is centered at \$58,000, which equals the
        population median. Is the sample median a biased or unbiased estimator? Explain.
        \writelines{2}
  \item For the estimator $\bar{x}$ from a population with $\sigma = 30$: how large
        must $n$ be so that $\sigma_{\bar{x}} \le 3$?
        \vspace{0.6in}
\end{enumerate}
\end{practicebox}

\end{document}
